\clearpage\section*{Correction EV4-A}\label{EV4-A-CORR}
\textbf{QCM :} 1 C, 2 C, 3 D, 4 C, 5 A, 6 A, 7 C, 8 B, 9 A, 10 A, 11 B, 12 A. Chaque réponse correcte : 0,5 point ; total 6.
\paragraph{Vérifications des automatismes.} Priorités ; somme de fractions ; soustraction des exposants ; coefficient multiplicateur 1,2 ; division par 0,8 ; produit 1,25 fois 0,8 ; résolution affine ; accroissement vertical divisé par l’horizontal ; moyenne ; angle divisé par 360 ; complément à 1 ; substitution entre parenthèses.
\subsection*{Exercice 1 -- Deux variables liées}
\begin{enumerate}
\item Les deux espérances valent 2. \textit{Barème : 1 point. Résultat et justification élémentaire attendus ; demi-point possible pour démarche correcte inachevée.}
\item $E(X^2)=\dfrac{14}{3}$, donc $V(X)=\dfrac{2}{3}$ ; même variance pour Y. X+Y=4, variance 0. \textit{Barème : 2 points. Méthode ou justification : 1 ; calcul et conclusion : 1.}
\item $0\ne \dfrac{2}{3}+\dfrac{2}{3}$. Le contre-exemple est sur un même univers ; Y dépend de X. \textit{Barème : 2 points. Méthode ou justification : 1 ; calcul et conclusion : 1.}
\end{enumerate}
\subsection*{Exercice 2 -- Comparer courbe et droite}
\begin{enumerate}
\item $h\prime=1-e^{-x}$. \textit{Barème : 1 point. Résultat et justification élémentaire attendus ; demi-point possible pour démarche correcte inachevée.}
\item Négative si x<0, positive si x>0 ; minimum h(0)=0. \textit{Barème : 2 points. Méthode ou justification : 1 ; calcul et conclusion : 1.}
\item Approximation 0,98 ; par défaut grâce à l’inégalité. \textit{Barème : 2 points. Méthode ou justification : 1 ; calcul et conclusion : 1.}
\end{enumerate}
\subsection*{Exercice 3 -- Vecteur normal et rayon}
\begin{enumerate}
\item $\dfrac{|3+4-17|}{5}=2$. \textit{Barème : 1 point. Résultat et justification élémentaire attendus ; demi-point possible pour démarche correcte inachevée.}
\item $\lambda =-\dfrac{10}{25}=-\dfrac{2}{5}$ ; H($\dfrac{11}{5}$;$\dfrac{13}{5}$). \textit{Barème : 1 point. Résultat et justification élémentaire attendus ; demi-point possible pour démarche correcte inachevée.}
\item $(x-1)^2+(y-1)^2=4$ ; pour H, $(\dfrac{6}{5})^2+(\dfrac{8}{5})^2=4$. \textit{Barème : 2 points. Méthode ou justification : 1 ; calcul et conclusion : 1.}
\end{enumerate}
\paragraph{Étalonnage.} Durée cible : 20–25 min pour les automatismes, 80–85 min pour les exercices, 10–15 min de relecture. Ces durées sont des estimations éditoriales à mesurer en classe. Accepter toute méthode correcte ; ne pas sanctionner deux fois une erreur numérique propagée avec cohérence. Le barème proposé ne vaut pas décision d’harmonisation de l’équipe.
